\section{Discussion}
\label{sec:discussion}

\subsection{Trade correction is part of the method}

Pakistan's net clinker exports are 12.7\% of the correctly period-tagged no-trade estimate
after conversion to \COtwo{}. Omitting them assigns calcination to the country where exported
clinker is ground rather than the country where it was produced. The corrected equation
therefore changes both the magnitude and the geographic boundary of 2.A.1. Stock movement
remains unresolved, but it is now visible as one term in the same physical balance.

The source audit also illustrates why period and source documentation are part of the
calculation. The reported 48.011 and 39.566 Mt values were consecutive observations in the
same official fiscal series, rather than competing estimates for one year. Correcting the
period changes the headline result more than the remaining parametric uncertainty and removes
a spurious structural choice.

\subsection{A complete sector starts with completeness accounting}

Pakistan's official results make the scope problem concrete. In 2022, chemical-industry
emissions were two-thirds as large as mineral-industry emissions. A cement calculation can
therefore be rigorous without representing the full IPPU sector. The framework avoids
conflating the two by assessing the evidence and method available for every 2.A--2.G category
before calculating a sector total.

This design also changes how public data are interpreted. A soda-ash production record is
useful for 2.B.7 only after the production route is known and must be distinguished from
soda-ash consumption under mineral uses. Total fertiliser output cannot estimate ammonia,
where process \COtwo{} recovery into urea is central. Billets identify steel output but not
the production route or carbon inputs. HFC consumption establishes chemical supply, whereas
annual emissions depend on equipment stocks and leakage. These distinctions preserve the
physical meaning of each category estimate.

\subsection{The CRT audit identifies immediate empirical priorities}

All tested CRT rows reproduce exactly from their reported activities, implied factors and
recovery terms. Its \SI{49.112}{\mega\tonne} result is consequently traceable to reported
activities, factors, recovery and omissions rather than a spreadsheet mistake. Three
empirical priorities follow.

First, the sixteen-company clinker census---complete for production since Fecto's exact
output was read from its annual report---should be completed with clinker purchases, sales
and stocks. The census already exceeds the low-clinker inference and slightly exceeds the
high case, so these data must reconcile reporting boundaries as well as replace Tier~1
activity. Second, plant ammonia and nitric-acid records can
reconstruct the official chemical-industry result while preserving recovered carbon and
abatement. Third, species- and application-level HFC imports, exports and equipment stocks can
support an estimate based on the Ozone consumption data. The small HFC value in the 2022 CRT
and the much larger consumption statistic have different accounting boundaries, but together
they make product-use completeness a high-priority audit. The CRT entry is smaller than
``small'' conveys and its shape matters: all of category~2.F is a single species, HFC-23,
entered at \SI{0.236955}{\tonne} as exactly \SI{1}{\percent} of a declared
\SI{23.69549}{\tonne} operating stock and filed under sub-category 2.F.1.a commercial
refrigeration, with every other species and sub-category \texttt{NE} or \texttt{NO}.
HFC-23 is an HCFC-22 by-product and a cascade low-temperature refrigerant rather than a
commercial refrigerant, so the species assignment is a reporting artefact, and Pakistan's
first Biennial Transparency Report then books 2.F, 2.C, 2.E and 2.G all at zero. The two
most recent official submissions therefore disagree about whether the category exists,
which is not a boundary difference.

The leak rate is the other decision that is easy to overlook, and
\texttt{outputs/pakemis\_fgas\_leak\_rate\_adjudication.csv} prices it. The
species-resolved bank applies one operating loss fraction per stratum---refrigeration
\SI{8}{\percent} per year, stationary air-conditioning \SI{6}{\percent}, mobile
air-conditioning \SI{15}{\percent}---while IPCC Table~7.9 publishes a rate for each
sub-application. Mobile air-conditioning maps one-to-one onto its table row and needs no
weighting. Stationary air-conditioning lumps chillers with residential and commercial
units, whose bands overlap, and \SI{6}{\percent} is admissible for either part.
Refrigeration lumps five sub-applications whose operation-emission bands span
\SIrange{0.1}{50}{\percent} per year with no common value, and the single
\SI{8}{\percent} lies outside two of them, from \num{16} to \num{80} times above
domestic refrigeration's whole band. For that part the rate is not uncertain but inadmissible, and
the split that would resolve it is not held.

What the envelope shows is that the rate decides timing rather than mass. Re-running the
bank across refrigeration's full band moves the operating channel from
\SI{0.013}{\mega\tonne} to \SI{2.062}{\mega\tonne} of \COtwo e, a factor of
one hundred and fifty-five, while the stratum total moves only
\SIrange{0.757}{2.192}{\mega\tonne}: gas the operating channel does not leak is still
emitted at retirement, because no recovery is credited anywhere. The deferred mass appears
instead as a closing bank of \SIrange{1.959}{12.559}{\mega\tonne} of \COtwo e. A low
leak rate is therefore not a smaller inventory but a larger future liability, and
category~2.F contributes zero to every reported total throughout.

\subsection{Facility geography is preferable to proxy geography}

The current provincial cement allocation is retained for compatibility with the national
inventory, but its earlier Punjab activity anchor is not period-matched. Cement kilns,
ammonia plants, nitric-acid plants, soda-ash works, steel mills and glass plants are point
sources. A shared facility inventory can assign process emissions to their physical sites and
link the same installation to Energy 1.A.2 combustion without mixing the two categories.
Diffuse product use, especially refrigeration and air-conditioning, can then be represented
through equipment stocks and population or service activity rather than an industrial point
proxy.

\subsection{The matched official comparison reveals a second cement discrepancy}

Calendar matching does not close the official cement row. Even the 0.95 public-data
configuration remains 19\% below the CRT, while the low-clinker model is 36\% below. Part of
the difference is visible in the activity data: the CRT uses 50.246 Mt against 44.331 Mt from
the monthly QIM reconstruction. The remainder is embedded in the CRT's 0.55539 t/t implied
factor and any clinker, stock or coverage treatment behind it. Because the workbook arithmetic
closes, resolving this difference requires the institutional derivation and company clinker
balance rather than further arithmetic verification.

\subsection{Transferability}

The framework is transferable because its accounting structure is independent of any one
national data source. Countries with incomplete industrial statistics commonly combine
national production totals, trade data, partial company reporting, ozone-substitute
consumption records and official inventory submissions. Higher-quality evidence can replace
aggregate activity assumptions without changing the material balances or category
definitions. Exact mass closure, category completeness and an auditable official crosswalk
remain valid as the quality of the underlying evidence improves.
